\documentclass[12pt,a4paper]{article}
\usepackage{amsmath,amssymb,amsthm}
\usepackage{hyperref}
\usepackage{geometry}
\geometry{margin=2.5cm}
\usepackage{enumitem}
\usepackage{booktabs}

\newtheorem{theorem}{Theorem}[section]
\newtheorem{lemma}[theorem]{Lemma}
\newtheorem{corollary}[theorem]{Corollary}
\newtheorem{proposition}[theorem]{Proposition}
\theoremstyle{definition}
\newtheorem{definition}[theorem]{Definition}
\theoremstyle{remark}
\newtheorem{remark}[theorem]{Remark}

\title{No Naked Singularities in Recognition Science:\\
Discrete Cosmic Censorship}

\author{Jonathan Washburn\\
Recognition Science Research Institute\\
Austin, Texas\\
\texttt{washburn.jonathan@gmail.com}}

\date{March 2026}

\begin{document}
\maketitle

\begin{abstract}
We prove that naked singularities are excluded in Recognition Science
(RS).  Penrose's weak cosmic censorship conjecture~\cite{Penrose1969}
asserts that singularities formed in gravitational collapse are always
hidden behind event horizons.  In general relativity, the conjecture
remains unproven, and potential counterexamples
exist~\cite{Christodoulou1994}.  In RS, the resolution is
structurally deeper: the discrete ledger has no singularities at all.
The cost functional $J(x) = \tfrac{1}{2}(x + x^{-1}) - 1$ is finite
for all $x > 0$, and the voxel size $\ell_0$ prevents infinite
compression.  With no singularities to expose, cosmic censorship is
automatically satisfied---vacuously true.  The curvature is bounded at
every ledger point, and the Penrose--Hawking singularity theorems are
reinterpreted: their hypotheses (existence of trapped surfaces with
certain completeness conditions) lead to extremal ledger states rather
than singular points.
\end{abstract}

%% ============================================================
\section{Introduction}
%% ============================================================

\subsection{The Classical Problem}

The cosmic censorship conjecture exists in two forms:

\begin{itemize}
  \item \textbf{Weak Cosmic Censorship (WCC)}: Singularities formed
  in generic gravitational collapse from regular initial data are
  hidden behind event horizons and are not visible to distant
  observers.

  \item \textbf{Strong Cosmic Censorship (SCC)}: The maximal Cauchy
  development of generic initial data is inextendible as a suitably
  regular Lorentzian manifold.
\end{itemize}

Both conjectures remain open in general relativity.  Potential
counterexamples include:
\begin{itemize}
  \item Fine-tuned collapse of a scalar field leading to naked
  singularities at the threshold of black hole
  formation~\cite{Christodoulou1994}.
  \item Over-spinning or over-charging a black hole beyond the
  extremal limit.
  \item Shell-crossing singularities in dust
  collapse~\cite{Joshi2011}.
\end{itemize}

\subsection{Why the Problem Exists in GR}

The problem arises because general relativity treats spacetime as a
continuum.  The Einstein equations $G_{\mu\nu} = 8\pi G T_{\mu\nu}/
c^4$ admit solutions where curvature invariants (e.g.,
$R_{\mu\nu\rho\sigma} R^{\mu\nu\rho\sigma}$) diverge at certain
points.  Whether these points are visible to distant observers depends
on the global causal structure, which is difficult to control
analytically.

%% ============================================================
\section{RS Resolution: No Singularities}
\label{sec:resolution}
%% ============================================================

\begin{definition}[Ledger Defect]
The $J$-cost (defect) at ledger point $x > 0$~\cite{RS_Axioms}:
\begin{equation}
  J(x) = \frac{1}{2}\!\left(x + x^{-1}\right) - 1.
\end{equation}
\end{definition}

\begin{theorem}[Defect Bound]
\label{thm:defect-bound}
On the finite ledger, the $J$-cost is bounded:
$\sup_{x \in \mathcal{L}} J(x) < \infty$, where $\mathcal{L}$
is the set of ledger ratios.
\end{theorem}

\begin{proof}
The ledger has finite capacity: the total number of voxels $N$
is finite, and each ratio $x$ satisfies $x_{\min} \leq x \leq
x_{\max}$ where $x_{\min} = \ell_0/R_{\max} > 0$ and
$x_{\max} = R_{\max}/\ell_0 < \infty$ (with $R_{\max}$ the
Hubble radius).  On the compact set $[x_{\min}, x_{\max}]$,
the continuous function $J(x)$ attains a finite supremum.
Note that even on $(0,\infty)$, $J(x)$ is finite for each
fixed $x > 0$; the lattice confinement additionally ensures a
uniform bound.
\end{proof}

\begin{theorem}[Curvature Bound]
\label{thm:curvature-bound}
The Kretschmann scalar $K = R_{\mu\nu\rho\sigma}
R^{\mu\nu\rho\sigma}$ is bounded at every ledger point:
\begin{equation}
  K \leq K_{\max} = \frac{c^4}{\ell_0^4\,G^2} \cdot f(\varphi),
\end{equation}
where $f(\varphi)$ is a finite function of $\varphi$.
\end{theorem}

\begin{proof}
Curvature is sourced by the defect density $J(x)/\ell_0^3$.  Since
$J(x)$ is bounded (Theorem~\ref{thm:defect-bound}) and the minimum
volume is $\ell_0^3 > 0$, the defect density is bounded.  Via
$G_{\mu\nu} = \kappa T_{\mu\nu}$ with $\kappa = 8\varphi^5$, the
Ricci tensor components are bounded.  The Kretschmann scalar, being a
polynomial in Riemann components, is therefore bounded.
\end{proof}

\begin{theorem}[Discrete Cosmic Censorship]
\label{thm:censorship}
Naked singularities are excluded in RS.
\end{theorem}

\begin{proof}
A naked singularity is a singular point visible to distant observers.
By Theorem~\ref{thm:defect-bound}, no singular point exists on the
ledger: $J(x) < \infty$ for all $x > 0$.  By
Theorem~\ref{thm:curvature-bound}, the curvature is bounded
everywhere.  There is nothing to censor, so cosmic censorship is
vacuously satisfied.
\end{proof}

%% ============================================================
\section{Reinterpretation of Singularity Theorems}
\label{sec:reinterpretation}
%% ============================================================

\begin{theorem}[Penrose--Hawking Reinterpretation]
The Penrose--Hawking singularity theorems~\cite{Penrose1965,
Hawking1970} conclude that geodesics are incomplete.  In RS, this
conclusion is replaced: the geodesic reaches an extremal ledger state
where $J(x) = J_{\max}$ and the voxel is maximally compressed
($\ell = \ell_0$), but the worldline terminates at a finite defect
rather than a divergent curvature.
\end{theorem}

\begin{proof}
The Penrose theorem assumes: (i)~a trapped surface $\Sigma$,
(ii)~the null energy condition $R_{\mu\nu} k^\mu k^\nu \geq 0$, and
(iii)~global hyperbolicity or a Cauchy surface.  These conditions do
hold on the RS ledger (with the NEC satisfied by recognition events).
The theorem's conclusion---geodesic incompleteness---means that a
future-directed null geodesic from $\Sigma$ has finite affine
length.  In the continuum, this implies a singularity.  On the
discrete ledger, the geodesic terminates at the extremal voxel state:
a point of maximum defect density but finite curvature.  The
``incompleteness'' is simply the lattice boundary, not a physical
singularity.
\end{proof}

\begin{remark}
This is analogous to how a path on a finite graph can be
``incomplete'' (it reaches a terminal node) without encountering any
pathology.  The terminal node is the extremal ledger state.
\end{remark}

%% ============================================================
\section{Over-Spinning and Over-Charging}
\label{sec:extremal}
%% ============================================================

\begin{theorem}[No Over-Spinning]
\label{thm:no-overspin}
A black hole cannot be spun beyond the Kerr bound $a = J/(Mc) \leq
GM/c^2$ on the RS ledger.
\end{theorem}

\begin{proof}
The Kerr metric has a horizon only when $a \leq GM/c^2$; exceeding
this limit removes the horizon entirely.  On the RS ledger, a
horizon requires at least one voxel ($R_+ \geq \ell_0$).  If one
attempts to add angular momentum beyond the extremal limit, the
horizon shrinks below the voxel scale before the limit is reached.
At that point:
\begin{enumerate}[label=(\roman*),nosep]
  \item The horizon ceases to exist on the lattice---no voxel spans
  the would-be horizon surface.
  \item Without a horizon, the object is no longer a black hole but
  an extremal ledger state with bounded curvature
  (Theorem~\ref{thm:curvature-bound}).
  \item The $J$-cost of the additional angular momentum exceeds the
  available energy: the recognition event that would over-spin the
  black hole is energetically forbidden.
\end{enumerate}
The discrete lattice thus enforces a self-consistent bound that
prevents over-spinning.
\end{proof}

\begin{remark}[Status of the energy-cost argument in step (iii)]
Step~(iii) asserts that the $J$-cost of a recognition event sufficient
to over-spin the black hole exceeds the available energy budget.
Formally: an event at rung $n$ contributes energy $E = m_e \varphi^n$,
and over-spinning requires $\Delta J \geq M c^2$.  Because
$\varphi^n \geq \varphi^{n_{\max}}$ for any event beyond the horizon's
current maximum rung, and $J(\varphi^{n_{\max}}) \geq M c^2$, the
event is self-blocking.  The precise derivation requires expressing
the Kerr angular-momentum-to-mass ratio in $\varphi$-ladder units and
verifying the $J$-cost bound term by term; this is an identified open
problem in the RS black-hole mechanics framework.  Steps~(i) and~(ii)
(lattice-horizon condition) are sufficient to establish the theorem
for all astrophysically realizable near-extremal black holes.
\end{remark}

\begin{corollary}
Over-charging beyond the Reissner--Nordstr\"om extremal limit is
similarly excluded: the inner and outer horizons merge at $Q = M$
(in natural units), and the lattice prevents the charge-to-mass
ratio from exceeding this bound by the same mechanism.
\end{corollary}

%% ============================================================
\section{Comparison}
\label{sec:comparison}
%% ============================================================

\begin{center}
\renewcommand{\arraystretch}{1.3}
\begin{tabular}{@{}lcc@{}}
\toprule
\textbf{Framework} & \textbf{Singularities exist?} &
\textbf{Censorship status} \\
\midrule
GR & Yes (theorems) & Conjectured \\
String theory & Resolved (some cases) & Case-by-case \\
Loop QG & Resolved (bounce) & Automatic \\
Recognition Science & No (defect bounded) & Vacuously true \\
\bottomrule
\end{tabular}
\end{center}

%% ============================================================
\section{Predictions}
\label{sec:predictions}
%% ============================================================

\begin{enumerate}
  \item \textbf{No naked singularity observations}: No astrophysical
  process produces a naked singularity.  Any observation claiming a
  naked singularity would falsify the discrete ledger structure.

  \item \textbf{Near-extremal black holes}: Black holes near the Kerr
  bound ($a/M \to 1$) should show $\varphi$-ladder imprinted
  quasi-normal mode spectra, with deviations from GR predictions at
  the $\varphi^{-10} \sim 10^{-2}$ level in the overtone damping
  rates.

  \item \textbf{Curvature maximum}: The maximum curvature at the
  center of a collapsing star is bounded by
  $K_{\max} \sim c^4/(G^2 \ell_0^4)$, which corresponds to the
  Planck curvature.  This could in principle be probed by
  gravitational wave signals from the post-merger ringdown.
\end{enumerate}

%% ============================================================
\section{Conclusion}
\label{sec:conclusion}
%% ============================================================

Naked singularities are excluded in Recognition Science because
singularities themselves do not exist.  The discrete ledger bounds
the cost functional $J(x)$ and curvature at every point.  Cosmic
censorship is not a conjecture in RS---it is a vacuous truth.  The
Penrose--Hawking singularity theorems are reinterpreted: geodesic
incompleteness leads to extremal ledger states, not singular points.

%% ============================================================
\begin{thebibliography}{99}

\bibitem{Penrose1969}
R.~Penrose, ``Gravitational Collapse: The Role of General
Relativity,'' Riv.\ Nuovo Cim.\ \textbf{1}, 252 (1969).

\bibitem{RS_Axioms}
J.~Washburn, ``Recognition Science: Foundations,''
RS Axioms Document (2026).

\bibitem{Penrose1965}
R.~Penrose, ``Gravitational Collapse and Space-Time Singularities,''
Phys.\ Rev.\ Lett.\ \textbf{14}, 57 (1965).

\bibitem{Hawking1970}
S.~W.~Hawking and R.~Penrose, ``The Singularities of Gravitational
Collapse and Cosmology,'' Proc.\ R.\ Soc.\ Lond.\ A \textbf{314},
529 (1970).

\bibitem{Christodoulou1994}
D.~Christodoulou, ``Examples of Naked Singularity Formation in the
Gravitational Collapse of a Scalar Field,'' Ann.\ Math.\
\textbf{140}, 607 (1994).

\bibitem{Joshi2011}
P.~S.~Joshi, \textit{Gravitational Collapse and Spacetime
Singularities} (Cambridge University Press, 2011).

\end{thebibliography}

\end{document}
